\chapterNote{Based on a talk given on February 10, 2013}

\chapter{The Delusion of Forever}
\tocChapterNote{Based on a talk given on February 10, 2013}
%\tocChapterNote{February 10, 2013}


Luang Por Pasanno's most recent teachings during this Winter Retreat have been exploring the five \emph{khandha}s, the five aggregates. The term aggregates refers to five categories of personal experience: form (the body), feeling, perception, volitional formations, and consciousness. These are the aggregates of existence.

The Buddha doesn't mince words in his teachings and with his straightforwardness and directness; he uses some pretty bold expressions, especially in reference to the khandhas. One of the phrases that often occurs in the teachings refers to the five aggregates as a disease, a dart, a tumor, a boil, a calamity. These are strong words. They point to the potential for the aggregates to cause suffering when we invest them with a sense of ownership and use them as a means for seeking satisfaction in the world.

\looseness=-1
One of the teachings from the Saṃyutta Nikāya is titled ``Yamaka'' \mbox{(SN 22.85)}, the name of one of the monks. To \mbox{summarize}, Yamaka is possessed of an annihilationist view and the monks try to dispossess him of this view. They are not successful and so they consult with Sāriputta, a chief disciple of the Buddha, who does an incredible job of pointing out his wrong view and directing him to understand that the Buddha's teaching is neither annihilationist nor eternalist.

Sāriputta then uses a simile that is very powerful. The story relates there is a wealthy householder, and there is a man who has something against him and wants to murder the wealthy man. He thinks to himself, ``Suppose I become his servant, get to know him, and work my way into his life.'' He approaches the wealthy man and offers his services as a servant. The householder accepts this offer, allows him to come into his life, and becomes familiar with him. Over time, with diligence, the servant attends to him, providing him with all he needs and gets to know the householder well. He rises in the morning before the householder, ready to greet him, and stays up until after he has gone to sleep. He is attentive and endearing in his speech and actions, slowly gaining the wealthy man's trust.

After some time, with a comfortable routine well-established, the householder lets his guard down completely. The servant then seizes the opportunity and stabs him in the back, killing him.

\enlargethispage{1\baselineskip}
Sāriputta then asks Yamaka, ``When this servant was initially offering his services, was he a murderer then, even though the householder didn't know it?'' Yamaka replies, ``Yes.''

``And when he was getting to know him, serving him, and setting the stage, was he a murderer then?'' ``Yes.''

``And when he finally stabbed him in the back when the householder wasn't looking, was he a murderer then?'' Yamaka answers, ``Yes, he was.'' So even though he didn't know it, there was a murderer in the house the entire time.

The murderer in this story is a metaphor for the five khandhas. We place a lot of trust in them, depend on them, and think they are here for us. They serve us well and provide us with pleasurable, warm experiences, and we become lulled by the pleasure they provide. All the while, we are not paying attention to the fact that at some point, they can and will turn against us. We willingly buy into this paradigm and believe in it. The khandhas are there for us when we go to bed. We sleep, dream, wake up and the khandhas are still there, ready to offer us what we \mbox{want to gratify our desires.}

We know that the teachings tell us not to put our faith in the aggregates, yet we still do. Time after time, we get pulled into the belief that we have a right to expect the khandhas to operate the way we want them to operate; that they should do what we want them to do. In the United States, we grow up with a constitutional right to life, liberty, and the pursuit of happiness, and we incorporate that into our belief system. And then something goes wrong.

Many of us are having physical health problems. Various muscles, joints, and internal organs are injured or deteriorating. The body doesn't obey us the way we think it should. The Buddha gave the monastics permission to use the four requisites, one of which is medicine for when we are sick. We have access to medical care, and we do what we can to \mbox{keep the body relatively healthy.}

The problem comes when we expect it to stay that way. We get a treatment, and maybe it doesn't work. We have an operation, and it fails. Perhaps the problem isn't what we thought it was, and it's going to take longer to fix than \mbox{we think}, or it's not going to get better at all.

One time, before I went in for surgery to repair a hernia, I thought, ``Oh, I'll be up and moving comfortably in a couple of days.'' For the surgery, I was given an anesthetic and morphine. I went to sleep and woke up a little while later. When I was more alert, the doctor came over to me and said, ``Well, it was bigger than we were anticipating, so you're going to be in recovery for a longer time.'' I had blissfully been put to sleep, woke up, and found out that I had been duped by this body.

\looseness=-1
It happens to all of us in different ways. We recover from an illness and we think, ``Now I'll be able to really get down to practice. My energy will return and everything will be fine.'' We believe the problem is cured and there is not going to be any more difficulty; we're home free. Maybe that assumption lasts for a little while, but then something else happens.

\enlargethispage{1\baselineskip}
It's not that we don't want to keep ourselves healthy and keep our bodies going as best we can. But if there is any expectation that it will be smooth sailing from then on, we are setting ourselves up for the next round of suffering. It is like putting a Band-Aid on a large, gaping wound, thinking it will heal it forever, or like plugging the hole in the Titanic with a piece of chewing gum. It's really a desperate fantasy.

Many years ago, a longtime friend of the monastery was hiking and camping in the mountains. One evening he lay down in a field to sleep and woke up a quadriplegic. He had been run over by a vehicle that was traveling across the field. He wasn't expecting that when he went to sleep that night. This friend has reflected that the rest of us with ``able bodies'' refer to people in his situation as ``disabled,'' but sometimes people we call ``disabled'' refer to us as TABs, ``temporarily able-bodied.'' That is a reminder that no condition is permanent.

\looseness=-1
Even the body itself inclines towards renewal, health, and revitalization, sometimes in the face of great odds. A strong image that remains in my mind is from when our neighbor was diagnosed with colon cancer and was starting to realize that it would be the cause of her demise. The cancer had spread to many places in her body, and she had surgery to remove part of the diseased area. At that time, she was clearly in the last few weeks of her life. I remember her lying there with her digestive system shutting down. She wasn't taking very much nourishment and was starting to waste away. Other organ systems were also shutting down. Even though she was getting IV fluids, the kidneys and circulatory system weren't functioning properly. The fluid was collecting in her legs and lungs, and at times she was having trouble breathing. We could see the body was packing it in.

\enlargethispage{1\baselineskip}
One time when the visiting nurse was changing the dressing on one of the surgical incisions from the operation, it was clear the wound was healing, it was getting better. Everybody was commenting, ``Oh, that looks really good. That's not going to cause you any more problems. We won't have to change the bandage anymore. It's healing up.'' It hit me as being so bizarre that she didn't have a lot longer to live---\mbox{her body's} major organ systems were all failing---but still, the body was working to heal itself, trying to mend this little wound. It struck me as so paradoxical.

The body itself seems hardwired to support us in this delusion of permanence. Then, when we turn our backs, it proves otherwise. It's the same with the other khandhas: feelings, perceptions, mental formations, and consciousness. We have an expectation, a hope, a fantasy that as soon \mbox{as we get} rid of an unpleasant feeling and find a way to replace it with a pleasant feeling, everything will be okay. So we search for something, either in the material realm or the mental realm, to provide a pleasant feeling. If we obtain it, we relax and think, ``Ah, this is the way it's supposed to be.'' This can even be the pleasantness of a ``good'' \mbox{meditation. We rate our meditations as a} ``bad'' meditation or a ``good'' meditation based on the level of bliss that we experience. We come to expect that a pleasant meditation is what we should have. Then, something changes and pleasant turns to unpleasant.

Regarding the aggregate of perception, we can cling to a perception so strongly that when something comes along to disprove it, it is very hard to change and perceive it differently. The Buddha refers to perceptions as mirages because it is so easy to view people and events through the lens of distorted impressions that are not accurate representations of those people and events.

Think of people you know. In terms of personal experience, people are just a perception in the mind---how easily that changes. Someone you have a lot of good times with can suddenly change and do something that takes you aback. Your perception changes dramatically and you experience him or her as a completely different person. All of a sudden, you find out that your faithful partner of twenty years is shacking up with the secretary or inviting the mailman in for lunch. What happened to that wonderful person in your life, who was always so warm and supportive and always there when you needed him or her?

More commonly, someone that you get along with has a bad day and says something sharp or angry, and all of a sudden he or she is an evil monster, at least for a period of time. Just when you think you have got it all worked out with the people around us and everybody is harmonizing, something happens and it changes. You're caught by surprise.

\looseness=-1
The fourth aggregate of mental formations includes our moods and emotions: how much control do we really have over them? We work with them and try to be more skillful by cultivating wholesome states of mind. That is part of the path. But when we think that they should be one way and that they shouldn't be another way, then we entangle ourselves even more. We create another mood or thought that is unpleasant, negative, or averse. We can slip into further despair in reaction to that negative mind state. Once again, we haven't lived up to our expectations, and we blame ourselves, compounding the mind state with another \mbox{layer of suffering.}

Consciousness, the fifth aggregate, is constantly changing. The nature of consciousness is to weave the illusion of continuity. The continuity of experience through consciousness gives us a sense of solidity and identity. Consciousness constantly pulls data from the six senses. For example, visual consciousness is ignited by an interesting object, but boredom quickly sets in. It is interesting for a flash and then the desire to experience the object fades just as quickly. An interesting sound arises, pleasant or unpleasant, and captivates our attention. We pull that into the mind and suddenly we are experiencing ear consciousness. The same is true with the other sense bases: smells, tastes, tactile sensations, and ideas of the mind. Attention is directed toward an object in one sense base, abandons it and immediately takes in another object, perhaps from a different sense base. We weave this together with the perception of ``this is who I am.'' All of these moments coming together in rapid succession create a sense of continuity, like a movie being projected onto a screen, and we believe in it.

This is how the five khandhas operate when they're affected by craving, clinging, and delusion. They keep us feeling secure, entertained, and feeling like this is who we are. We constantly renew this perception over and over again, but unexpectedly, the khandhas betray us.

\looseness=-1
However, there is also an interesting paradox in the teachings. In order to free ourselves from this confusion, the teachings point us back to the five khandhas. They are not only our ``murderer,'' but also the way out of this whole \mbox{mass of suffering.}

There is a sutta where a young deva of stunning beauty appears in front of the Buddha, illuminating Jeta's Grove \mbox{(SN 2.26)}. His name is Rohitassa; he pays respects to the Buddha and asks him a question. He says, ``Bhante, is it possible by traveling that one can reach the end of the world where there is no birth, no aging, no death, no passing away, and no rebirth?'' The Buddha replies, ``No, Rohitassa, there is no way by traveling that one can reach that end of the world where there is no birth, aging, death, or rebirth.'' Rohitassa says, ``Oh, it is wonderful that you say that because in a past life I was a seer, a sage, referred to as a `skywalker.' With one stride, I could step across an ocean, and my speed was rapid like that of a light, swift arrow soaring through the sky. I had a lifespan of a hundred years, and I spent my life coursing through the cosmos, with my quick and long stride, trying to find the end of the world where there is no birth, aging, death, and rebirth. But I died before I got to the end of the world.'' The Buddha responds, ``Yes, there is no way that by traveling you can reach the end of the world, the end of suffering. I tell you that it's just in this fathom-long body where there is the world, the origin of the world, the cessation of the world, and the way leading to the \mbox{cessation of the world.''}

The path to the end of suffering is found by exploring the experience within our bodies. The examination of these five khandhas, of what they really are and are not, is the way leading to the end of suffering.

\enlargethispage{1\baselineskip}
We have to be willing to throw our whole life into our spiritual practice, every bit of energy we have. We have to be willing to make that commitment because the habits are strong. The tendency to want to find satisfaction through the experience of these five khandhas is very strong. At the same time, the full-on commitment has to be made with compassion, kindness, and gentleness. Commitment is not something that can be forced on us. It has to be taken up with much care, attention, and clarity.

It is through investigation of the khandhas that we learn to know what has been fooling us all along. We don't just disregard the need for this clear and complete examination. We don't run away from it. We turn around and look carefully at the khandhas. We turn around and face the murderer, which is not easy to do. We look at all that we have been investing in for so long, with fresh eyes. We need to bring a lot of mindfulness and alertness to the process.

At the same time, we have to be willing to recondition our perceptions. We can cultivate skillful mind states that result in more useful perceptions. The new perceptions are also subject to change and deterioration, but we use them in a skillful way. We use mindfulness and alertness to develop all the path factors, including renunciation, kindness, and gentleness. We use the aggregates that have been fooling us, turn them into instruments for a deeper understanding, and see them in their true nature. We see through the illusion of permanence and solidity and realize that the aggregates cannot be a refuge. We can, however, use them in a skillful way in order to understand how they really operate.

\enlargethispage{1\baselineskip}
When we accurately see how the khandhas operate, over and over again, disenchantment with searching for any permanent sense of satisfaction from the aggregates arises, and identification with them weakens. We have avoided seeing the aggregates clearly for so long and now we have the opportunity to look at them squarely in the face. We use them for our long-term happiness and ultimate liberation.
